\documentclass[journal]{IEEEtran}
\usepackage[utf8]{inputenc}
\usepackage{amsmath,amssymb}
\usepackage{graphicx}
\usepackage{booktabs}
\usepackage{hyperref}
\usepackage{cite}
\usepackage{algorithmic}
\usepackage{algorithm}

\begin{document}

\title{Your IEEE Journal Paper Title}

\author{
  First~Author,~\IEEEmembership{Member,~IEEE,}
  Second~Author,~\IEEEmembership{Senior~Member,~IEEE}
  \thanks{F. Author is with the Department of ECE, University One, City, Country (e-mail: first@example.com).}
  \thanks{S. Author is with the Department of CS, University Two, City, Country.}
}

\markboth{IEEE Transactions on XXX, Vol.~XX, No.~X, Month~Year}
{Author \MakeLowercase{\textit{et al.}}: Paper Title}

\maketitle

\begin{abstract}
This paper presents ... Provide a concise abstract of 150-250 words covering the problem, approach, key results, and conclusions.
\end{abstract}

\begin{IEEEkeywords}
Keyword one, keyword two, keyword three, keyword four.
\end{IEEEkeywords}

\section{Introduction}
\IEEEPARstart{T}{his} paper addresses the problem of ...

\section{Related Work}
Review relevant literature.

\section{System Model}
Describe the system model and problem formulation.

\section{Proposed Method}
\subsection{Overview}
Provide a high-level description.

\subsection{Algorithm}
Detail the proposed algorithm.

\section{Experimental Results}
\subsection{Setup}
Describe evaluation methodology.

\subsection{Performance Comparison}
Present comparative results.

\section{Conclusion}
Summarize the paper and discuss future work.

\bibliographystyle{IEEEtran}
\bibliography{references}

\end{document}
